%% kapa: exam=2 round=36 year=2025 period=3 subject=compensation_law subject_ko=감정평가 및 보상법규 seq=1 number=1 points=40
%% source: Q-Net 감정평가사 기출문제 (cst00309) artlSeq=5239325 file=raw/1008/5239325/2025년 제36회 감정평가사 2차 감정평가 및 보상법규.pdf
%% conversion: kordoc 4.24.51 -> tools/gen2.py (auto); review=proofread-gpt-sol; flags=none
\begin{longquestion}{【문제 1】}
주식회사 甲은 30년 넘게 서울특별시 A구에서 콘크리트 제조 공장을 허가 받아 적법하게 운영하고 있다. 甲이 창고로 사용하고 있는 부지(이하 ‘이 사건 부지’라 함)는 백제 한성기의 왕성 유적으로 알려진 국가지정문화유산 B유적지의 보호구역에 있다. 학계 등은 B유적지의 완전한 복원을 위해서는 이 사건 부지의 수용 필요성을 오래전부터 제기해 왔고, 이 사건 부지에 있는 창고의 증축과정에서 백제 토성의 유물이 발견되기도 하였다. 한편, 관련 법령에 따라 B유적지를 관리하는 A구청장 乙은 B유적지의 정비 및 복원을 위하여 이 사건 부지에 대한 공익사업을 계획하고 국토교통부장관에게 사업인정 신청을 하였다. 국토교통부장관은 「공익사업을 위한 토지 등의 취득 및 보상에 관한 법률」(이하 ‘토지보상법’이라 함)에 따라 사업인정(이하 ‘이 사건 사업인정’이라 함)을 하고 이를 고시하였다. 甲은 다음과 같은 이유로 이 사건 사업인정 및 고시가 위법하다고 주장한다. \par
1. 이 사건 부지가 수용되면 공장 운영의 차질 및 콘크리트 공급 감소로 인한 건설 현장의 대혼란을 야기할 수 있으므로 문화유산 보호의 이유로는 수용될 수 없다. \par
2. 乙은 국가지정문화유산을 관리하는 행정청이지 수용 권한이 있는 것은 아니다. 다음 물음에 답하시오. (40점) \par
\medskip\noindent 물음 1) 토지보상법상 다음의 법적 성격을 설명하시오. (10점) \par
① 사업인정 전과 후의 ‘보상에 관한 협의’ \par
② 사업인정 후 ‘협의취득’ \par
\medskip\noindent 물음 2) 甲 주장(1. 및 2.)의 타당성 여부를 검토하시오. (20점) \par
\medskip\noindent 물음 3) 이 사건 부지에 대한 수용이 진행될 경우, 甲의 영업손실에 대한 보상 요건을 설명하시오. (10점) \par
\begin{framedtext}
참조 조문 \par
<공익사업을 위한 토지 등의 취득 및 보상에 관한 법률> \par
제16조(협의) 사업시행자는 토지등에 대한 보상에 관하여 토지소유자 및 관계인과 성실하게 협의하여야 하며, 협의의 절차 및 방법 등 협의에 필요한 사항은 대통령령으로 정한다. \par
제20조(사업인정) ① 사업시행자는 제19조에 따라 토지등을 수용하거나 사용하려면 대통령령으로 정하는 바에 따라 국토교통부장관의 사업인정을 받아야 한다. \par
<공익사업을 위한 토지 등의 취득 및 보상에 관한 법률 [별표]> \par
1. 법 제20조에 따라 사업인정을 받아야 하는 공익사업 \par
\medskip\noindent (1) 「공간정보의 구축 및 관리 등에 관한 법률」에 따른 기본측량의 실시 \par
\medskip\noindent (2) 「공공토지의 비축에 관한 법률」에 따라 한국토지주택공사가 공공개발용 토지의 비축사업계획을 승인받은 공공개발용 토지의 취득 \par
\medskip\noindent (3) 「국립대학법인 서울대학교 설립ㆍ운영에 관한 법률」에 따른 국립대학법인 서울대학교의 학교 용지 확보 \par
\medskip\noindent (4) 「국립대학법인 인천대학교 설립ㆍ운영에 관한 법률」에 따른 국립대학법인 인천대학교의 학교 용지 확보 \par
\medskip\noindent (5) 「규제자유특구 및 지역특화발전특구에 관한 규제특례법」에 따른 특화사업 \par
\medskip\noindent (6) 「농어업재해대책법」에 따른 응급조치 \par
\medskip\noindent (7) 「대기환경보전법」 제4조에 따라 고시된 측정망설치계획에 따른 환경부장관 또는 시ㆍ도지사의 측정망 설치 \par
\medskip\noindent (8) 「문화유산의 보존 및 활용에 관한 법률」, 「자연유산의 보존 및 활용에 관한 법률」에 따른 문화유산과 자연유산의 보존ㆍ관리 \par
<공익사업을 위한 토지 등의 취득 및 보상에 관한 법률 시행규칙> \par
제45조(영업손실의 보상대상인 영업) 법 제77조제1항에 따라 영업손실을 보상하여야 하는 영업은 다음 각 호 모두에 해당하는 영업으로 한다. \par
1. 사업인정고시일등 전부터 적법한 장소(무허가건축물등, 불법형질변경토지, 그 밖에 다른 법령에서 물건을 쌓아놓는 행위가 금지되는 장소가 아닌 곳을 말한다)에서 인적ㆍ물적시설을 갖추고 계속적으로 행하고 있는 영업. 다만, 무허가건축물등에서 임차인이 영업하는 경우에는 그 임차인이 사업인정고시일등 1년 이전부터 「부가가치세법」 제8조에 따른 사업자등록을 하고 행하고 있는 영업을 말한다. \par
2. 영업을 행함에 있어서 관계법령에 의한 허가등을 필요로 하는 경우에는 사업인정고시일등 전에 허가등을 받아 그 내용대로 행하고 있는 영업 \par
<문화유산의 보존 및 활용에 관한 법률> \par
제83조(토지의 수용 또는 사용) ① 국가유산청장이나 지방자치단체의 장은 문화유산의 보존ㆍ관리를 위하여 필요하면 지정문화유산이나 그 보호구역에 있는 토지, 건물, 나무, 대나무, 그 밖의 공작물을 「공익사업을 위한 토지 등의 취득 및 보상에 관한 법률」에 따라 수용(收用)하거나 사용할 수 있다. \par
\end{framedtext}
\end{longquestion}
